\section{Scope, Method, and Qualifications}
\label{app:scope}

\subsection{Purpose, Audience, and Claim Boundary}

This Catholic single-passage scriptural study examines John 20:3--10 alongside Luke 24:12---controlled by John 20:1--2 and New Testament burial and textile parallels (Matt 27:59--66, 28:11--15; Mark 14:51--52, 15:46, 16:1--8; Luke 4:17--20, 19:20, 23:53--56, 24:1--14, 21--24, 34; John 11:43--44, 19:38--42; Acts 10:11, 11:5, 19:12; Rom 6:9)---for believers, students, and teachers seeking primary linguistic, historical, liturgical, and patristic evidence. It is a study aid, not an official liturgical book, critical edition, catechism, or magisterial act, and carries no ecclesiastical approbation or imprimatur.

Two claims stand under test: what the Greek and Latin texts establish about the two Apostles, the \latin{othonia}, the \latin{soudarion}, and ``he saw and believed''; and whether the popular modern account of a first-century Jewish dining-napkin signal has any ancient foundation. Forensic, radiocarbon, textile, and provenance debates concerning the Shroud of Turin and the Sudarium of Oviedo are excluded as belonging to distinct material-relic disciplines; omitting them neither affirms nor denies their authenticity.

\subsection{Editions, Languages, and Translation Practice}

Controlling English Scripture quotations follow the public-domain Douay--Rheims Bible (Challoner revision) verbatim; comparative references cite the King James Version (1611), the Revised Version (1895), and the World English Bible (Catholic Edition) by name. Latin Scripture follows the Clementine Vulgate. Greek texts are checked in the Robinson--Pierpont Byzantine Textform (RP2018 Unicode CSV with inline Nestle--Aland divergences) and the SBL Greek New Testament (Holmes v1.2, for Luke); New Testament lemma counts are verified against MorphGNT SBLGNT v6.12; Luke 24:12's textual status is checked in Westcott and Hort (1882 Appendix, p. 71); and lexicography is checked in Thayer's corrected edition.

Greek and Hebrew terms appear in scholarly Latin-script transliteration with macrons (\latin{\={e}}, \latin{\={o}}) where helpful so both print and web surfaces render without glyph substitution. Historical public-domain English translations are quoted directly for Chrysostom (Marriott, NPNF I.14), Augustine (Gibb and Innes, NPNF I.7; Salmond, NPNF I.6), Cyril of Alexandria (Randell, \emph{Library of Fathers} vol. 2), Aquinas's \emph{Catena Aurea} (Oxford 1874), Cornelius a Lapide (Mossman et al., 1908), the Roman Catechism (Donovan, 1829), the Babylonian Talmud (Cohen, 1921), and Suetonius (Rolfe, Loeb). Plain literal English renderings of Latin sources (Ambrose, Gregory the Great, Bede, Theophylact's Latin column, Aquinas's \emph{Super Ioannem}) are supplied directly from the checked Latin texts. Edward Caswall's 1849 English translation of \emph{Victimae paschali laudes} (\emph{Lyra Catholica}, pp. 233--234) is reproduced as a historical literary rendering, not an official liturgical text; the Latin follows the 1962 \emph{Missale Romanum}.

\subsection{Bounded Negative Findings, Source Ceilings, and Rights}

Every negative historical finding regarding the folded-napkin custom is bounded strictly to the corpora inspected and remains open to correction by any verified ancient evidence. Five source ceilings apply: (1) electronic Latin texts of Augustine (\texttt{augustinus.it}), Ambrose, Gregory, and Bede (Migne PL via \emph{Corpus Corporum}), and Aquinas (\emph{Corpus Thomisticum}, Turin 1952 base) are working transcriptions rather than \emph{Corpus Christianorum} or Leonine critical editions; (2) Theophylact on John 20 was read in Migne's PG 124 facing Latin column (OCR with four checked character repairs: \latin{habuissent}, \latin{deponendi}, \latin{bitumine}, \latin{insumpsisset}) rather than the uncollated Greek column; (3) Lightfoot's four-volume 1859 Oxford edition was examined in OCR where an isolated recognition miss cannot be excluded in principle, though the intact John 20 lemmata in volume 3 pass directly from Ver. 5 to Ver. 12 with no lemma on Ver. 7; (4) \emph{m. Berakhot} 8:3 (Torat Emet Hebrew), \emph{m. Sanhedrin} 6:1 (Sefaria Community CC0 alongside Torat Emet Hebrew), and \emph{b. Berakhot} 42a and 52b (Cohen 1921) were read in dated Sefaria public deliveries rather than printed volumes; and (5) the 7 October 2006 date on David Bivin's \emph{Jerusalem Perspective} page comes from site metadata inside a 7 March 2013 Internet Archive capture, the 21 May 2007 \emph{RecipesAndMore} forward is a circulation lower bound rather than proof of origin, and \emph{TruthOrFiction} carries both an inline update of 28 January 2008 and a byline of 17 March 2015.

All quoted historical Scripture translations, Greek/Latin texts, patristic and scholastic translations, and nineteenth-century reference works are public domain (or CC0 / CC BY-SA as noted in References). Copyrighted modern web pages are referenced analytically without bulk verbatim quotation. No external specialist peer review or ecclesiastical review has been performed.
